\chapter{Les interactions fondamentales}\label{ch:g11:fundamental-interactions}

Frottez un ballon sur vos cheveux et il ramasse des bouts de papier ---
battant, avec quelques centimètres carrés de caoutchouc, l'attraction
gravitationnelle de la planète entière. Cette victoire facile est la
raison pour laquelle les \omterm{def:g11:fundamental-interactions:ladder}{atomes} tiennent
ensemble et pour laquelle les noyaux lourds finissent par éclater
(\cref{ch:g11:nucleus-radioactivity}). L'univers tourne sur exactement
quatre interactions~; ce chapitre les rencontre, les pèse l'une contre
l'autre, et assigne à chacune son étage.

\section{La matière des quarks aux galaxies}

\begin{definition}[L'échelle de structure]\label{def:g11:fundamental-interactions:ladder}
La matière est construite en étages, chacun assemblé à partir de celui
du dessous~: trois \emph{quarks}\index{quark} (ponctuels, plus petits
que \qty{e-18}{m}) se lient en un \emph{nucléon}\index{nucléon} ---
proton ou neutron, d'environ \qty{e-15}{m} de diamètre~; les nucléons
s'empilent en un \emph{noyau}\index{noyau} (\qty{e-15}{m} à
\qty{e-14}{m})~; un noyau plus son nuage d'électrons est un
\emph{atome}\index{atome} (\qty{e-10}{m})~; les atomes se lient en
\emph{molécules}\index{molécule} et cristaux (à partir de \qty{e-9}{m}),
qui s'assemblent en objets du quotidien (\qty{1}{m}), planètes
(\qty{e7}{m}), systèmes planétaires (\qty{e11}{m}), galaxies
(\qty{e21}{m}) et l'univers observable (\qty{e26}{m}).
\end{definition}

\begin{definition}[Les quatre interactions]\label{def:g11:fundamental-interactions:four}
Une \emph{interaction fondamentale}\index{interaction fondamentale} est
une \omterm{def:g10:inertia:force}{force} non réductible à d'autres
\omterm{def:g10:inertia:force}{forces}~; il y en a exactement quatre~:
la \emph{gravitation}\index{gravitation}, l'\emph{interaction
électromagnétique}\index{interaction électromagnétique},
l'\omterm{def:g11:fundamental-interactions:strong}{interaction forte} et
l'\omterm{def:g11:fundamental-interactions:weak}{interaction faible}.
\end{definition}

\begin{remark}\label{rem:g11:fundamental-interactions:empty}
Les étages sont séparés par des vides~: un
\omterm{def:g11:fundamental-interactions:ladder}{atome} est $10^{5}$ fois plus large que son
\omterm{def:g11:fundamental-interactions:ladder}{noyau} --- ramenez le
\omterm{def:g11:fundamental-interactions:ladder}{noyau} à une bille de \qty{1}{cm} et le
nuage électronique fait un kilomètre de diamètre. Quelque chose doit
tenir ces étages clairsemés, et la gravité n'y est presque jamais pour
rien.
\end{remark}

\begin{omfigure}
\begin{tikzpicture}[font=\small, x=0.24cm]
  \draw[-Stealth] (-1,0) -- (46.5,0);
  \foreach \x/\l/\d in {0/{10^{-18}}/2pt, 3/{10^{-15}}/13pt,
                        8/{10^{-10}}/2pt, 18/{10^{0}}/2pt,
                        25/{10^{7}}/2pt, 29/{10^{11}}/13pt,
                        39/{10^{21}}/2pt, 44/{10^{26}}/13pt}
    \draw (\x,0.1) -- (\x,-0.1) node[below=\d] {$\l$};
  \foreach \x/\n in {0/quark, 3/noyau, 8/atome, 18/humain, 25/Terre,
                     29/{Soleil--Terre}, 39/galaxie, 44/univers}{
    \fill[omDef] (\x,0) circle (1.6pt);
    \node[rotate=45, anchor=west, inner sep=1.5pt] at (\x,0.24) {\n};}
\end{tikzpicture}

{\small L'échelle de structure sur un axe en puissances de dix~:
quarante-quatre décades séparent le
\omterm{def:g11:fundamental-interactions:ladder}{quark} de l'univers observable. Tailles en
\omterm{def:g10:orders-of-magnitude:unit}{mètres}.}
\end{omfigure}

\section{Charge électrique}

\begin{definition}[Charge électrique]\label{def:g11:fundamental-interactions:charge}
La \emph{charge électrique}\index{charge électrique} est la propriété de
la matière sur laquelle agit l'\omterm{def:g11:fundamental-interactions:four}{interaction
électromagnétique}, comme la masse est la propriété sur laquelle agit
la \omterm{def:g11:fundamental-interactions:four}{gravitation}. Elle se mesure en
\emph{coulombs}\index{coulomb} (\unit{C}) et se présente sous
\emph{deux signes}, positif et négatif, qui s'annulent par addition~;
les charges de même signe se repoussent, de signes opposés
s'attirent.
\end{definition}

\begin{definition}[Charge élémentaire]\label{def:g11:fundamental-interactions:elementary}
La charge est granulaire~: toute charge observée est un multiple entier
de la \emph{charge élémentaire}\index{charge élémentaire}
$e = \qty{1.60e-19}{C}$. Le proton porte $+e$, l'électron $-e$, le
neutron $0$~; une charge $q$ contient $N = \abs{q}/e$ charges
élémentaires.
\end{definition}

\begin{proposition}[Conservation de la charge]\label{prop:g11:fundamental-interactions:conservation}
La charge totale d'un système isolé ne change jamais~: la charge n'est
ni créée ni détruite, seulement transférée --- en pratique par les
électrons, les porteurs les plus légers.
\end{proposition}
\admitted

\begin{remark}\label{rem:g11:fundamental-interactions:conservation}
Aucune violation n'a jamais été observée~; la raison profonde, une
symétrie de l'électromagnétisme, est donnée dans les volumes
universitaires.
\end{remark}

\begin{definition}[Électrisation par frottement et par contact]\label{def:g11:fundamental-interactions:charging}
Frotter deux isolants arrache des électrons de l'un vers l'autre, laissant
des charges opposées d'égale grandeur (\emph{électrisation par
frottement}\index{électrisation par frottement})~; toucher un objet
chargé à un objet neutre partage la charge (\emph{électrisation par
contact}\index{électrisation par contact}). Seuls les électrons se
déplacent~; la charge totale est conservée.
\end{definition}

\begin{example}[Compter les électrons]\label{ex:g11:fundamental-interactions:balloon}
Un ballon frotté sur des cheveux acquiert $q = \qty{-1.6e-8}{C}$~: il a
\emph{gagné} $N = \num{1.6e-8} / \num{1.6e-19} = \num{1.0e11}$
électrons, et les cheveux restent avec exactement $+\qty{1.6e-8}{C}$.
\end{example}

\section{Loi de Coulomb}

\begin{theorem}[Loi de Coulomb]\label{thm:g11:fundamental-interactions:coulomb}
Deux charges ponctuelles $q_1$ et $q_2$ distantes de $d$ exercent l'une
sur l'autre des \omterm{def:g10:inertia:force}{forces} le long de la droite qui
les joint, d'égales grandeurs et de directions opposées ---
répulsives pour des signes identiques, attractives pour des signes
opposés --- avec
\[
F = k\,\frac{\abs{q_1 q_2}}{d^{2}},
\qquad k = \qty{8.99e9}{N.m^2/C^2}.
\]
\end{theorem}
\admitted

\begin{remark}\label{rem:g11:fundamental-interactions:torsion}
Coulomb a établi la loi en 1785 en mesurant la torsion d'une fibre
\omterm{def:g11:lenses-and-eye:lens}{mince}~; pourquoi l'\omterm{def:g10:orders-of-magnitude:scinot}{exposant}
est exactement $2$ est répondu dans les volumes universitaires.
\end{remark}

\begin{omfigure}
\begin{tikzpicture}[font=\small]
  \draw[omDef, thick, fill=omDef!20] (0,0) circle (0.3) node[black] {$+$};
  \draw[omDef, thick, fill=omDef!20] (2.6,0) circle (0.3) node[black] {$+$};
  \draw[-Stealth, omThm, very thick] (-0.3,0) -- (-1.5,0)
    node[midway, above] {$\vect F_{2/1}$};
  \draw[-Stealth, omThm, very thick] (2.9,0) -- (4.1,0)
    node[midway, above] {$\vect F_{1/2}$};
\end{tikzpicture}\qquad
\begin{tikzpicture}[font=\small]
  \draw[omDef, thick, fill=omDef!20] (0,0) circle (0.3) node[black] {$+$};
  \draw[omProp, thick, fill=omProp!25] (2.6,0) circle (0.3) node[black] {$-$};
  \draw[-Stealth, omThm, very thick] (0.3,0) -- (1.1,0)
    node[midway, above] {$\vect F_{2/1}$};
  \draw[-Stealth, omThm, very thick] (2.3,0) -- (1.5,0)
    node[midway, above] {$\vect F_{1/2}$};
\end{tikzpicture}

{\small Charges de même signe se repoussent (à gauche), de signes
opposés s'attirent (à droite)~; les deux
\omterm{def:g10:inertia:force}{forces} ont toujours des grandeurs égales et des
directions opposées.}
\end{omfigure}

\begin{remark}[Un jumeau formel de la gravitation]\label{rem:g11:fundamental-interactions:twin}
La loi de Coulomb est, symbole pour symbole, la loi de
\omterm{def:g11:fundamental-interactions:four}{gravitation} universelle
(\cref{ch:g10:universal-gravitation}) avec des charges à la place des
masses~:
\begin{center}
\begin{tabular}{lll}
source & masse $m$ & charge $q$ \\
loi & $F = G\,\dfrac{m_1 m_2}{d^2}$ & $F = k\,\dfrac{\abs{q_1 q_2}}{d^2}$ \\[1.5ex]
constante & $G = \qty{6.67e-11}{N.m^2/kg^2}$ & $k = \qty{8.99e9}{N.m^2/C^2}$ \\
signe & toujours attractive & attractive ou répulsive \\
\end{tabular}
\end{center}
Les deux différences mènent tout ce qui suit~: l'électricité est
écrasamment plus forte, et elle seule peut s'annuler elle-même.
\end{remark}

\begin{example}[Deux protons, les deux forces à la fois]\label{ex:g11:fundamental-interactions:protons}
Deux protons ($m_p = \qty{1.67e-27}{kg}$) sont distants de
$d = \qty{1.0e-15}{m}$ --- voisins nucléaires. Répulsion électrique~:
\[
F_E = \num{8.99e9} \times
\frac{(\num{1.60e-19})^{2}}{(\num{1.0e-15})^{2}}
\approx \qty{2.3e2}{N}
\]
--- le \omterm{def:g10:universal-gravitation:weight}{poids} d'une valise de $23$~kg, porté
par une particule de $10^{-27}$
\omterm{def:g10:orders-of-magnitude:unit}{kilogrammes}. Attraction gravitationnelle~:
$F_G = \num{6.67e-11} \times (\num{1.67e-27})^{2} /
(\num{1.0e-15})^{2} \approx \qty{1.9e-34}{N}$. Le rapport
$F_E / F_G = k e^{2} / (G m_p^{2}) \approx \num{1.2e36}$ est
indépendant de la distance ($d^{2}$ s'annule)~: entre particules
élémentaires, la gravité est $10^{36}$ fois trop faible pour compter ---
à \emph{n'importe quelle} distance.
\end{example}

\section{Les interactions forte et faible}

L'exemple laisse un scandale~: un
\omterm{def:g11:fundamental-interactions:ladder}{noyau} d'hélium tient deux protons qui se
repoussent avec des dizaines de \omterm{def:g10:inertia:force}{newtons} --- et
pourtant l'hélium existe.

\begin{definition}[L'interaction forte]\label{def:g11:fundamental-interactions:strong}
L'\emph{interaction forte}\index{interaction forte} lie les
\omterm{def:g11:fundamental-interactions:ladder}{quarks} en
\omterm{def:g11:fundamental-interactions:ladder}{nucléons} et les
\omterm{def:g11:fundamental-interactions:ladder}{nucléons} en noyaux. À \qty{e-15}{m} elle
est attractive et bien plus forte que la répulsion électrique (des
milliers de \omterm{def:g10:inertia:force}{newtons} entre
\omterm{def:g11:fundamental-interactions:ladder}{nucléons})~; au-delà de quelques $10^{-15}$
\omterm{def:g10:orders-of-magnitude:unit}{mètres} elle s'éteint~: sa
\emph{portée}\index{portée (d'une interaction)} est d'environ
\qty{e-15}{m}. Elle saisit protons et neutrons indifféremment et ignore
la charge.
\end{definition}

\begin{omfigure}
\begin{tikzpicture}[font=\small]
  \draw[omDef, thick, fill=omDef!20] (0,0) circle (0.34) node[black] {p};
  \draw[omDef, thick, fill=omDef!20] (3.2,0) circle (0.34) node[black] {p};
  \draw[-Stealth, omThm, thick] (-0.34,0) -- (-1.7,0)
    node[midway, above] {$F_E \approx \qty{230}{N}$};
  \draw[-Stealth, omThm, thick] (3.54,0) -- (4.9,0) node[midway, above] {$F_E$};
  \draw[-Stealth, omMeth, very thick] (0.2,0.55) -- (1.45,0.55)
    node[midway, above] {forte};
  \draw[-Stealth, omMeth, very thick] (3.0,0.55) -- (1.75,0.55)
    node[midway, above] {forte};
  \draw[Stealth-Stealth] (0,-0.65) -- (3.2,-0.65)
    node[midway, below] {$d = \qty{1.0e-15}{m}$};
\end{tikzpicture}

{\small Le bras de fer à l'intérieur d'un
\omterm{def:g11:fundamental-interactions:ladder}{noyau}~: à $10^{-15}$~m l'attraction forte
(des milliers de \omterm{def:g10:inertia:force}{newtons}) bat la répulsion
électrique.}
\end{omfigure}

\begin{remark}[Pourquoi les noyaux existent --- et pourquoi les lourds cèdent]\label{rem:g11:fundamental-interactions:nuclei}
Les deux \omterm{def:g10:inertia:force}{forces} dans le
\omterm{def:g11:fundamental-interactions:ladder}{noyau} ne scalent pas de la même façon. La
colle forte est à courte portée~: un
\omterm{def:g11:fundamental-interactions:ladder}{nucléon} ne lie que ses quelques voisins
dans \qty{e-15}{m}. La répulsion électrique est à longue portée~: chaque
proton pousse sur \emph{tous les autres} protons. À mesure que les
noyaux grossissent, la répulsion s'accumule plus vite que la colle~: les
noyaux lourds ont besoin de neutrons supplémentaires (colle sans
charge), et au-delà d'environ $80$ protons même cela échoue --- les
noyaux les plus lourds vivent à crédit
(\cref{ch:g11:nucleus-radioactivity}).
\end{remark}

\begin{definition}[L'interaction faible]\label{def:g11:fundamental-interactions:weak}
L'\emph{interaction faible}\index{interaction faible}, de
\omterm{def:g11:fundamental-interactions:strong}{portée} encore plus courte, laisse un
neutron se transformer en proton (et inversement)~: elle pilote la
radioactivité $\beta$ du \cref{ch:g11:nucleus-radioactivity}.
\end{definition}

\section{Qui règne à quel étage}

\begin{method}[Auditer une échelle]\label{met:g11:fundamental-interactions:audit}
Pour trouver quelle interaction règne sur une structure de taille $d$,
vérifier la \omterm{def:g11:fundamental-interactions:strong}{portée} et l'annulation~:
\begin{enumerate}
  \item $d \leq \qty{e-14}{m}$~: l'\emph{\omterm{def:g11:fundamental-interactions:strong}{interaction
        forte}} règne --- elle bat la répulsion électrique, et la
        gravité est hors jeu d'un facteur $10^{36}$.
  \item des \omterm{def:g11:fundamental-interactions:ladder}{atomes} aux montagnes
        ($\qty{e-10}{m}$ à $\qty{e4}{m}$)~: la
        \omterm{def:g10:inertia:force}{force} forte est hors de
        \omterm{def:g11:fundamental-interactions:strong}{portée}~; l'\emph{\omterm{def:g11:fundamental-interactions:four}{interaction
        électromagnétique}} règne --- liaisons, cohésion,
        \omterm{def:g10:inertia:inventory}{frottement}, contact.
  \item planètes et au-delà ($d \geq \qty{e6}{m}$)~: la matière est
        neutre à une précision fantastique, donc l'électricité
        s'annule~; la masse n'a qu'un signe et s'ajoute toujours --- la
        \emph{\omterm{def:g11:fundamental-interactions:four}{gravitation}} règne sur le ciel.
\end{enumerate}
L'\omterm{def:g11:fundamental-interactions:weak}{interaction faible} ne tient aucun étage~:
elle ne lie rien, mais arbitre des transformations (désintégration
$\beta$).
\end{method}

\begin{omfigure}
\begin{tikzpicture}[font=\small, x=0.26cm]
  \draw[-Stealth] (-1,0) -- (41.5,0);
  \foreach \x/\l in {3/{10^{-15}}, 8/{10^{-10}}, 18/{10^{0}},
                     24/{10^{6}}, 39/{10^{21}}}
    \draw (\x,0.1) -- (\x,-0.1) node[below] {$\l$};
  \draw[omProp, fill=omProp!35] (3,0.3) rectangle (4,0.75);
  \node[anchor=west, omProp] at (4.4,0.52) {forte};
  \draw[omDef, fill=omDef!25] (8,0.95) rectangle (22,1.4);
  \node[omDef] at (15,1.17) {électromagnétique};
  \draw[omThm, fill=omThm!25] (24,1.6) rectangle (39,2.05);
  \node[omThm] at (31.5,1.82) {gravitation};
\end{tikzpicture}

{\small Qui règne où~: forte pour les noyaux, électromagnétique des
\omterm{def:g11:fundamental-interactions:ladder}{atomes} aux montagnes,
\omterm{def:g11:fundamental-interactions:four}{gravitation} pour les planètes et au-delà.
Tailles en \omterm{def:g10:orders-of-magnitude:unit}{mètres}.}
\end{omfigure}

\section{Exercices}

\begin{exercise}[$\star$]\label{exo:g11:fundamental-interactions:1}
Donner l'\omterm{def:g10:orders-of-magnitude:oom}{ordre de grandeur} de chaque taille et
son étage de l'échelle~: proton \qty{1.7e-15}{m}~;
\omterm{def:g11:fundamental-interactions:ladder}{molécule} d'eau \qty{2.8e-10}{m}~; grain de
sable \qty{2e-4}{m}~; Terre \qty{1.3e7}{m}~; Voie lactée \qty{9.5e20}{m}.
\end{exercise}

\begin{exercise}[$\star$]\label{exo:g11:fundamental-interactions:2}
Un ballon frotté sur des cheveux porte $q = \qty{-4.8e-9}{C}$. A-t-il
gagné ou perdu des électrons, et combien~? Quelle est la charge des
cheveux ensuite, et quelle loi le dit~?
\end{exercise}

\begin{exercise}[$\star$]\label{exo:g11:fundamental-interactions:3}
Des charges $q_1 = \qty{+2.0e-6}{C}$ et $q_2 = \qty{-3.0e-6}{C}$ sont
distantes de \qty{30}{cm}. Calculer la
\omterm{def:g10:inertia:force}{force}~; attractive ou répulsive~?
Comparer les \omterm{def:g10:inertia:force}{forces} sur $q_1$ et sur $q_2$.
\end{exercise}

\begin{exercise}[$\star$]\label{exo:g11:fundamental-interactions:4}
Deux charges s'attirent avec une \omterm{def:g10:inertia:force}{force} $F_0$.
Que devient-elle si~: la distance double~; la distance est divisée par
$3$~; une charge double~; les deux charges \emph{et} la distance
doublent~?
\end{exercise}

\begin{exercise}[$\star$]\label{exo:g11:fundamental-interactions:5}
Nommer l'interaction responsable de~: (a)~la Lune orbitant autour de la
Terre~; (b)~la cohésion d'un cristal de sel~; (c)~la cohésion d'un
\omterm{def:g11:fundamental-interactions:ladder}{noyau} d'hélium~; (d)~la désintégration
$\beta$ du carbone-14~; (e)~le ballon de
l'\cref{exo:g11:fundamental-interactions:2} collé au mur.
\end{exercise}

\begin{exercise}[$\star\star$]\label{exo:g11:fundamental-interactions:6}
Les deux protons d'un \omterm{def:g11:fundamental-interactions:ladder}{noyau} d'hélium sont
distants de $d = \qty{2.0e-15}{m}$ ($m_p = \qty{1.67e-27}{kg}$). Calculer
leur répulsion électrique, leur attraction gravitationnelle, et le
rapport. Qu'est-ce qui maintient le
\omterm{def:g11:fundamental-interactions:ladder}{noyau} entier~?
\end{exercise}

\begin{exercise}[$\star\star$]\label{exo:g11:fundamental-interactions:7}
Dans un \omterm{def:g11:fundamental-interactions:ladder}{atome} d'hydrogène, l'électron
($m_e = \qty{9.11e-31}{kg}$) est à $d = \qty{5.3e-11}{m}$ du proton
($m_p = \qty{1.67e-27}{kg}$). Calculer les
\omterm{def:g10:universal-gravitation:force}{forces} électrique et gravitationnelle entre
eux et leur rapport. Qu'est-ce qui tient les
\omterm{def:g11:fundamental-interactions:ladder}{atomes} ensemble~?
\end{exercise}

\begin{exercise}[$\star\star$]\label{exo:g11:fundamental-interactions:8}
La sphère A porte $q_A = \qty{+8.0e-9}{C}$~; une sphère B identique est
neutre. On les met en contact, puis on les sépare de \qty{10}{cm}.
Quelle charge porte chacune (vérifier la conservation de la charge)~?
Calculer la \omterm{def:g10:inertia:force}{force} entre elles --- attractive
ou répulsive~?
\end{exercise}

\begin{exercise}[$\star\star$]\label{exo:g11:fundamental-interactions:9}
Deux charges identiques $q = \qty{1.0e-6}{C}$ se repoussent avec
\qty{0.90}{N}~: à quelle distance sont-elles~? Où la
\omterm{def:g10:inertia:force}{force} tombe-t-elle à \qty{0.10}{N}~?
\end{exercise}

\begin{exercise}[$\star\star$]\label{exo:g11:fundamental-interactions:10}
Deux protons dans des \omterm{def:g11:fundamental-interactions:ladder}{molécules} voisines
sont distants de $d = \qty{1.0e-10}{m}$. Calculer leur répulsion
électrique. Que contribue l'\omterm{def:g11:fundamental-interactions:strong}{interaction
forte} à cette distance~? Quelle interaction lie les
\omterm{def:g11:fundamental-interactions:ladder}{molécules}~?
\end{exercise}

\begin{exercise}[$\star\star$]\label{exo:g11:fundamental-interactions:11}
Chacune de deux pièces de monnaie détient environ $\num{e23}$
électrons. Déplacer un électron sur un \emph{million} entre elles et
placer les pièces à \qty{1.0}{m}. Calculer la charge de chaque pièce,
puis la \omterm{def:g10:inertia:force}{force}. Le
\omterm{def:g10:universal-gravitation:weight}{poids} de quelle masse l'égale-t-il~?
Conclure sur la neutralité de la matière ordinaire.
\end{exercise}

\begin{exercise}[$\star\star\star$]\label{exo:g11:fundamental-interactions:12}
Un \omterm{def:g11:fundamental-interactions:ladder}{noyau} d'uranium a un diamètre
\qty{1.4e-14}{m} et $92$ protons. Calculer la répulsion entre deux
protons aux extrémités opposées. L'\omterm{def:g11:fundamental-interactions:strong}{interaction
forte} peut-elle lier cette paire --- pourquoi non~? Combien de paires
de protons se repoussent-elles~? Pourquoi les noyaux lourds ont-ils
besoin de neutrons supplémentaires, et pourquoi au-delà d'une certaine
taille même les neutrons ne peuvent-ils les sauver~?
\end{exercise}

\begin{exercise}[$\star\star\star$]\label{exo:g11:fundamental-interactions:13}
Deux balles de masse \qty{1.0}{g} pendent à des fils de \qty{50}{cm}
attachés au même point. Dotées de la même charge $q$, elles s'écartent
de \qty{6.0}{cm}. Chaque balle est en équilibre sous son
\omterm{def:g10:universal-gravitation:weight}{poids}, la \omterm{def:g10:inertia:inventory}{tension}
du fil et la \omterm{def:g10:inertia:force}{force} électrique~: montrer que
$F = mg\tan\theta$ ($\theta$~: angle du fil à la verticale) et qu'ici
$\tan\theta \approx 0.060$. Calculer $F$, puis $q$, puis le nombre de
\omterm{def:g11:fundamental-interactions:elementary}{charges élémentaires} déplacées.
\end{exercise}

\begin{exercise}[$\star\star\star$]\label{exo:g11:fundamental-interactions:14}
Données~: $M_E = \qty{5.97e24}{kg}$, $M_M = \qty{7.35e22}{kg}$, distance
Terre--Lune $d = \qty{3.84e8}{m}$. Calculer la
\omterm{def:g10:universal-gravitation:force}{force gravitationnelle} Terre--Lune. La
gravité est coupée~: quelles charges $+Q$ sur la Terre et $-Q$ sur la
Lune maintiennent la même attraction~? Combien d'électrons vaut $Q$, et
combien pèsent-ils ($m_e = \qty{9.11e-31}{kg}$)~? Commenter.
\end{exercise}

\begin{exercise}[$\star\star\star$]\label{exo:g11:fundamental-interactions:15}
Dans un cristal de sel, des ions Na$^{+}$ et Cl$^{-}$ (charges $\pm e$)
sont distants de $d = \qty{2.8e-10}{m}$~; un ion Cl$^{-}$ a une masse
\qty{5.9e-26}{kg}. Calculer la \omterm{def:g10:inertia:force}{force} entre
ions voisins et le \omterm{def:g10:universal-gravitation:weight}{poids} d'un ion Cl$^{-}$~;
donner le rapport. Puis expliquer comment la gravité gagne aux grandes
échelles --- aucun sommet de montagne à \qty{e4}{m} --- bien que chaque
liaison la batte de quinze ordres de grandeur.
\end{exercise}

\section{Problème~: L'audit des quatre forces}

\begin{problem}[{Problème du week-end --- pourquoi le monde ne s'effondre
ni ne s'envole~: les quatre interactions auditées étage par étage,
jusqu'à la neutralité de la matière à deux parties près sur $10^{18}$}]
\label{pb:g11:fundamental-interactions:1}
Les \omterm{def:g11:fundamental-interactions:ladder}{atomes} n'implosent pas, les noyaux
tiennent pour la plupart, la Lune ni ne s'écrase ni ne s'échappe~: ce
problème audite à qui en revient le crédit. Données~:
$e = \qty{1.60e-19}{C}$,
$k = \qty{8.99e9}{N.m^2/C^2}$, $G = \qty{6.67e-11}{N.m^2/kg^2}$,
$m_p = \qty{1.67e-27}{kg}$, $m_e = \qty{9.11e-31}{kg}$,
$M_E = \qty{5.97e24}{kg}$, $M_M = \qty{7.35e22}{kg}$, distance
Terre--Lune \qty{3.84e8}{m}.

\medskip
\noindent\textbf{Partie I --- L'échelle.}
\begin{enumerate}
  \item Lister les étages de l'échelle du
        \omterm{def:g11:fundamental-interactions:ladder}{quark} à la galaxie, un
        \omterm{def:g10:orders-of-magnitude:oom}{ordre de grandeur} de taille chacun.
  \item Modèle réduit~: le \omterm{def:g11:fundamental-interactions:ladder}{noyau} devient
        une bille de \qty{1}{cm}~; quelle largeur a l'
        \omterm{def:g11:fundamental-interactions:ladder}{atome}~?
  \item Quelle fraction du volume d'un
        \omterm{def:g11:fundamental-interactions:ladder}{atome} son
        \omterm{def:g11:fundamental-interactions:ladder}{noyau} occupe-t-il~?
  \item Un cheveu fait \qty{70}{\micro m} d'épaisseur. Combien
        d'\omterm{def:g11:fundamental-interactions:ladder}{atomes} le traversent~?
  \item Combien de \omterm{def:g10:energy-conservation:power}{puissances} de dix du plafond
        du \omterm{def:g11:fundamental-interactions:ladder}{quark} (\qty{e-18}{m}) à une
        galaxie (\qty{e21}{m})~?
\end{enumerate}

\medskip
\noindent\textbf{Partie II --- L'étage électrique.}
\begin{enumerate}[resume]
  \item Dans un \omterm{def:g11:fundamental-interactions:ladder}{atome} d'hydrogène l'électron
        est à $d = \qty{5.3e-11}{m}$ du proton~: calculer la
        \omterm{def:g10:inertia:force}{force} électrique entre eux.
  \item Calculer la \omterm{def:g10:universal-gravitation:force}{force gravitationnelle}
        entre eux, et le rapport des deux. Laquelle tient l'
        \omterm{def:g11:fundamental-interactions:ladder}{atome} ensemble~?
  \item Que change-t-on si la gravité est coupée à l'intérieur des
        \omterm{def:g11:fundamental-interactions:ladder}{atomes}~? Et si l'électricité
        l'est~?
  \item Pourquoi les \omterm{def:g10:inertia:force}{forces} de «~contact~» de la
        vie quotidienne --- le sol qui pousse sous vos pieds, le
        \omterm{def:g10:inertia:inventory}{frottement} --- sont-elles
        électromagnétiques déguisées~?
  \item Chacune de deux pièces détient environ $\num{e23}$ électrons.
        Transférer un électron sur un \emph{milliard} entre elles, les
        placer à \qty{1.0}{m}~: calculer les charges et la
        \omterm{def:g10:inertia:force}{force}. Conclusion~?
  \item Quelle interaction règne sur les étages de \qty{e-10}{m} à
        \qty{e4}{m}, et pourquoi l'\omterm{def:g11:fundamental-interactions:strong}{interaction
        forte} ne concurrence-t-elle pas~?
\end{enumerate}

\medskip
\noindent\textbf{Partie III --- L'étage nucléaire.}
\begin{enumerate}[resume]
  \item Deux protons sont distants de \qty{1.0e-15}{m} dans un
        \omterm{def:g11:fundamental-interactions:ladder}{noyau}~: calculer leur répulsion
        électrique.
  \item Sans frein, quelle accélération cette
        \omterm{def:g10:inertia:force}{force} donnerait-elle à un proton~?
        Comparer avec $g$.
  \item Lister les trois propriétés dont l'
        \omterm{def:g11:fundamental-interactions:strong}{interaction forte} a besoin pour
        expliquer pourquoi les noyaux existent, pourquoi les neutrons
        sont saisis aussi, et pourquoi les noyaux sont minuscules.
  \item Dans un \omterm{def:g11:fundamental-interactions:ladder}{noyau} d'uranium (diamètre
        \qty{1.4e-14}{m}, $92$ protons), calculer la répulsion entre
        deux protons aux extrémités opposées. L'
        \omterm{def:g11:fundamental-interactions:strong}{interaction forte} peut-elle lier
        cette paire directement~? Expliquer pourquoi la répulsion gagne
        quand les noyaux grossissent, et à quoi servent les neutrons.
  \item En une phrase, comme promis~: que fait l'
        \omterm{def:g11:fundamental-interactions:weak}{interaction faible}~?
\end{enumerate}

\medskip
\noindent\textbf{Partie IV --- L'étage astronomique, et le verdict.}
\begin{enumerate}[resume]
  \item Calculer la \omterm{def:g10:universal-gravitation:force}{force gravitationnelle}
        entre la Terre et la Lune.
  \item La Terre contient environ $\num{1.8e51}$ protons (et autant
        d'électrons), la Lune environ $\num{2.2e49}$. Quelle fraction
        $f$ de la charge protonique de chaque corps, laissée
        non annulée, ferait égaler la
        \omterm{def:g10:inertia:force}{force} électrique à la gravitationnelle~?
  \item Pourquoi la \omterm{def:g11:fundamental-interactions:four}{gravitation}, le faible
        de la question~7, règne-t-elle en astronomie~? Deux raisons.
  \item Verdict --- une phrase par étage
        (\omterm{def:g11:fundamental-interactions:ladder}{noyau},
        \omterm{def:g11:fundamental-interactions:ladder}{atome} à montagne, planète et
        au-delà)~: qu'est-ce qui le tient ensemble~? Puis répondre au
        titre~: pourquoi le monde ni ne s'effondre ni ne s'envole~?
\end{enumerate}
\end{problem}
